\section{Introduction}\label{sec:introduction}
\subsection{Problem and motivation}
Lecture material is organised for a lecture: slide bullets, notes, code, formulas and recordings that assume a teacher in the room. A microlearning module needs something different, one narrow objective with an explanation, an example and questions aligned to it. Writing such modules by hand is repetitive, and a module drafted by a general language model may contain statements that the course material does not support.

GENAILMS (Generative Artificial Intelligence for Learning Management Systems) is a tool for lecturers. A lecturer uploads course documents, asks for a module, checks each statement against the passage it cites, edits the result and publishes it to Moodle as a Book and quizzes. Students reach the published activities in Moodle and can ask questions about the course material through the same tool.

The project reduces authoring effort while keeping every generated statement traceable to the lecturer's own material and leaving the final decision with the lecturer.
\subsection{Objectives}
\begin{enumerate}[label=O\arabic*.,leftmargin=*,nosep]
\item Build the complete workflow from course document to published Moodle activity.
\item Generate modules from passages retrieved in the lecturer's documents, with an inline citation to the exact passage for each important statement.
\item Let the lecturer inspect sources and edit the module while the first AI version is kept.
\item Integrate with Moodle through its identity and capability checks and through LTI 1.3.
\item Deploy as stateless services that can scale on AWS within a \$200 credit.
\item Measure retrieval quality, answer quality, latency, resource use and scalability, and define a blinded evaluation of module quality with human and model raters.
\end{enumerate}
\subsection{Module structure}\label{sec:module-structure}
A module targets one terminal learning objective (LO) with enabling sub-objectives. The cognitive level follows the revised Bloom taxonomy\citep{bloom}, and the parts are aligned with the objective following constructive alignment\citep{alignment}. A module has seven parts:
\begin{enumerate}[leftmargin=*,nosep]
\item \textbf{Learning Objective}: the expected behaviour and its Bloom level.
\item \textbf{Content}: a short lesson with one heading per sub-objective.
\item \textbf{Slides}: the lesson as a presentation, exportable to PPTX.
\item \textbf{Example}: a worked illustration of the concept.
\item \textbf{Practice}: questions with immediate feedback.
\item \textbf{Quiz}: assessment questions with feedback after submission.
\item \textbf{Explanation}: the reasoning behind each answer.
\end{enumerate}
\subsection{Scope}\label{sec:scope}
The evaluation corpus is the five subjects of \code{docs/documents}: Advanced Programming, Calculus 1, Machine Learning, Operating Systems and Scientific Socialism, in PDF and PPTX slides and lecture recordings. It replaced the four computing subjects of the midterm report after the advisor's feedback. The tool accepts PDF, DOCX, PPTX and HTML documents and common audio and video formats, which are transcribed with timestamps. Chapter detection in textbooks and learning-outcome studies with students are out of scope.
\subsection{Contributions}
\begin{itemize}[leftmargin=*,nosep]
\item A working authoring tool, from upload through generation, review, editing and Moodle publication, with LTI 1.3 entry for lecturers and students and roster synchronisation (O1, O3, O4).
\item Source-conditioned generation with physical citations: each citation resolves to a stored chunk and its page, slide or timestamp (O2).
\item A deployment of stateless services described in Terraform for AWS, with load and failure tests that located and removed the bottlenecks on one machine (O5).
\item Measurements on the five-subject corpus: retrieval (Hit Rate@$K$, MRR, nDCG), course chat faithfulness and relevancy with Ragas, a first pilot of the module evaluation, latency, throughput and failure behaviour (O6).
\end{itemize}
\subsection{Report structure}
Section~\ref{sec:background} reviews related work and the measurements behind the main technical choices. Sections~\ref{sec:requirements} and~\ref{sec:usecases} give requirements and use cases. Sections~\ref{sec:architecture} to~\ref{sec:ui} describe the architecture, data model, pipelines, Moodle integration and security, API and interface. Section~\ref{sec:evaluation} defines the evaluation method, Section~\ref{sec:results} reports the results, and Section~\ref{sec:conclusion} concludes and lists remaining work.
